\section{System Architecture}

The system is a modular pipeline; each stage is an independent Python module so components (e.g.\ the local optimiser) can be swapped without
touching the rest.

\begin{enumerate}
    \item \textbf{Map generation} (\texttt{utils/map\_generator.py}) --- produces a random
    binary occupancy grid with free start/goal cells.
    \item \textbf{Probabilistic map} (\texttt{src/prob\_map.py}) --- converts the grid into a
    continuous Gaussian risk field and derives an inflated planning grid.
    \item \textbf{Global planner} (\texttt{src/a\_star.py}) --- A* search on the inflated grid
    returns a waypoint path.
    \item \textbf{Local planner} (\texttt{src/mppi.py}, primary; \texttt{src/dwa.py} as a
    baseline) --- at each control step, evaluates candidate motions against a cost function
    and returns a velocity command.
    \item \textbf{Navigation loop} (\texttt{src/navigation.py}) --- drives the robot with a
    pure-pursuit look-ahead target along the path, integrating the motion model until the
    goal is reached.
    \item \textbf{Visualisation} (\texttt{utils/visualisation.py}) --- renders the grid, risk
    field, global path and executed trajectory.
\end{enumerate}

MPPI is the operational planner; the \texttt{PLANNER} flag in \texttt{main.py} can switch to
the DWA baseline for comparison. Both expose the same call signature, so the navigation loop
is agnostic to which one runs.
